%%% FIELD proof-published-class-calendar-loading
Put \(\delta_n=n^{-q/2}\), write \(r_i=r_i(W_i)\), and write \(M=M(\Pi)\). By the overlap condition in the published setup, \(r_i=0\) off \(\mathcal S\) and
\[
0\leq r_i\leq c_*^{-1}.
\]
The weighted outcome terms also have uniformly bounded second moments. Indeed, consistency, conditioning on \(\mathscr Y_i\), and \(\pi_i(s)\geq c_*\) give, for every \(t\),
\[
\begin{aligned}
\mathbb E\!\left[r_i^2Y_{it}^2\right]
&=\sum_{s\in\mathcal S}\mathbb E\!\left[
  \mathbf 1\{W_i=s\}\frac{\Pi_s^2}{\pi_i(s)^2}Y_{it}(s)^2
\right]\\
&=\sum_{s\in\mathcal S}\Pi_s^2
  \mathbb E\!\left[\frac{Y_{it}(s)^2}{\pi_i(s)}\right]
\leq c_*^{-1}M_2\sum_{s\in\mathcal S}\Pi_s^2
\leq c_*^{-1}M_2.
\end{aligned}
\]
Because \(T\) and \(p\) are fixed and every entry of \(D_s\) and \(J\) is bounded, every scalar entry needed below from
\[
r_i,\quad r_iD_{W_i},\quad r_iD_{W_i}^{\top}JD_{W_i},
\quad r_iJY_i,\quad r_iD_{W_i}^{\top}JY_i
\]
therefore has a second moment bounded by a constant independent of \(i\) and \(n\).

Let \(U_i\) be any one of these scalar entries. By the definition of the maximal-correlation coefficient,
\[
\left|\operatorname{Cov}(U_i,U_j)\right|
\leq \rho_{ij}
\sqrt{\operatorname{Var}(U_i)\operatorname{Var}(U_j)}
\leq C\rho_{ij}.
\]
Consequently the published weak-dependence condition implies
\[
\operatorname{Var}\!\left\{n^{-1}\sum_{i=1}^{n}
(U_i-\mathbb E U_i)\right\}
\leq Cn^{-2}\sum_{i=1}^{n}\sum_{j=1}^{n}\rho_{ij}
=O(n^{-q})=O(\delta_n^2).
\]
For every \(A>0\), the elementary inequality
\(\mathbf 1\{|Z|>A\delta_n\}\leq Z^2/(A^2\delta_n^2)\)
therefore yields
\[
\mathbb P_n\!\left(
\left|n^{-1}\sum_i(U_i-\mathbb E U_i)\right|>A\delta_n
\right)\leq \frac{C}{A^2}.
\]
Thus every required scalar sample moment equals its expectation plus \(O_P(\delta_n)\). The same statement holds for the finitely many entries of every required vector and matrix.

For every supported path \(s\) and every integrable function \(f(\mathscr Y_i)\), the definition of the known propensity gives the exact change-of-law identity
\[
\begin{aligned}
\mathbb E\!\left[r_i\mathbf 1\{W_i=s\}f(\mathscr Y_i)\right]
&=\mathbb E\!\left[
  \mathbb E\!\left[
   \left.\mathbf 1\{W_i=s\}\frac{\Pi_s}{\pi_i(s)}f(\mathscr Y_i)
   \right|\mathscr Y_i\right]
\right]\\
&=\Pi_s\mathbb E[f(\mathscr Y_i)].
\end{aligned}
\]
In particular, with
\[
s_n=n^{-1}\sum_{i=1}^{n}r_i,
\qquad
S_{D,n}=n^{-1}\sum_{i=1}^{n}r_iD_{W_i},
\]
the preceding moment bound gives
\[
s_n=1+O_P(\delta_n),
\qquad
S_{D,n}=M+O_P(\delta_n).
\]
Since \(s_n\) is bounded away from zero with probability tending to one,
\[
\overline D_n:=\frac{\sum_i r_iD_{W_i}}{\sum_i r_i}
=\frac{S_{D,n}}{s_n}=M+O_P(\delta_n).
\]

The grand intercept together with \(\sum_i\alpha_i=0\) and \(\sum_t\lambda_t=0\) spans all unit constants and common time effects. Eliminating these nuisance coefficients first applies \(J\) within each unit and then removes the weighted cross-unit mean. Hence, whenever the resulting normal matrix is nonsingular, the weighted least-squares normal equations give
\[
\widehat\tau^{\mathrm{ex}}=\widehat G_n^{-1}\widehat g_n,
\]
where
\[
\widehat G_n=\frac{1}{nT}\sum_{i=1}^{n}r_i
(D_{W_i}-\overline D_n)^{\top}J(D_{W_i}-\overline D_n)
\]
and
\[
\widehat g_n=\frac{1}{nT}\sum_{i=1}^{n}r_i
(D_{W_i}-\overline D_n)^{\top}JY_i.
\]
The weighted outcome mean disappears from the latter equation because
\(\sum_i r_i(D_{W_i}-\overline D_n)=0\). Expanding the first display and using the exact weighted mean gives
\[
\widehat G_n
=\frac{1}{nT}\sum_i r_iD_{W_i}^{\top}JD_{W_i}
  -\frac{s_n}{T}\overline D_n^{\top}J\overline D_n.
\]
The change-of-law identity, the moment rate, and
\(s_n=1+O_P(\delta_n)\), \(\overline D_n=M+O_P(\delta_n)\) now imply
\[
\begin{aligned}
\widehat G_n
&=\frac1T\sum_{s\in\mathcal S}\Pi_sD_s^{\top}JD_s
  -\frac1T M^{\top}JM+O_P(\delta_n)\\
&=\frac1T\sum_{s\in\mathcal S}\Pi_s
  (D_s-M)^{\top}J(D_s-M)+O_P(\delta_n)\\
&=G(\Pi)+O_P(\delta_n).
\end{aligned}
\]

For the numerator, add and subtract \(M\):
\[
\widehat g_n
=\frac{1}{nT}\sum_i r_i(D_{W_i}-M)^{\top}JY_i
 -(\overline D_n-M)^{\top}
   \left\{\frac{1}{nT}\sum_i r_iJY_i\right\}.
\]
The vector in braces is \(O_P(1)\): its centered part is \(O_P(\delta_n)\), and its expectation is uniformly bounded by the displayed second-moment bound and Cauchy--Schwarz. The second term is therefore \(O_P(\delta_n)\). For the first term, the change-of-law identity and the stated unconditional mean model give
\[
\begin{aligned}
&\mathbb E\!\left[\frac1T r_i(D_{W_i}-M)^{\top}JY_i\right]\\
&\quad=\frac1T\sum_{s\in\mathcal S}\Pi_s(D_s-M)^{\top}J
\left\{\mu_i+\sum_{t=1}^{T}\sum_{\ell=0}^{p}
e_t\tau_{i,t,\ell}(D_s)_{t,\ell}\right\}\\
&\quad=\sum_{t=1}^{T}\sum_{\ell=0}^{p}
h_{t,\ell}(\Pi)\tau_{i,t,\ell},
\end{aligned}
\]
where \(\mu_i=(\mu_{i1},\ldots,\mu_{iT})^{\top}\). The baseline term vanishes because
\(\sum_s\Pi_s(D_s-M)=0\), and the last equality is the definition of \(h_{t,\ell}(\Pi)\). Averaging over \(i\) and applying the moment rate proves
\[
\widehat g_n=b_n+O_P(\delta_n),
\qquad
b_n:=\sum_{t=1}^{T}\sum_{\ell=0}^{p}
h_{t,\ell}(\Pi)\overline\tau_{t,\ell,n}.
\]

The deterministic moving target \(b_n\) is uniformly bounded under the stated conditions. To see this without an additional effect-size assumption, let \(w^{(0)}\) be the all-zero path. The second-moment bound and Cauchy--Schwarz imply
\[
|\mu_{it}|=|\mathbb E Y_{it}(w^{(0)})|\leq M_2^{1/2}.
\]
For every legal pair \((t,\ell)\), meaning \(1\leq t-\ell\leq T\), let \(w^{(t,\ell)}\) have its sole nonzero coordinate at \(t-\ell\). The mean model then gives
\[
\tau_{i,t,\ell}
=\mathbb E Y_{it}(w^{(t,\ell)})-\mathbb E Y_{it}(w^{(0)}),
\qquad
|\tau_{i,t,\ell}|\leq 2M_2^{1/2}.
\]
For an illegal pair, \((D_s)_{t,\ell}=0\) for every \(s\), so \(h_{t,\ell}(\Pi)=0\). Since there are finitely many legal pairs and \(h_{t,\ell}(\Pi)\) is fixed,
\[
\sup_n\|b_n\|
\leq 2M_2^{1/2}\sum_{\substack{1\leq t\leq T,\ 0\leq\ell\leq p\\1\leq t-\ell\leq T}}
\|h_{t,\ell}(\Pi)\|<\infty.
\]

Let \(\lambda_{\min}(G)>0\) be the smallest eigenvalue of \(G=G(\Pi)\). Since the dimension \(d\) is fixed and \(\widehat G_n-G=O_P(\delta_n)=o_P(1)\), with probability tending to one
\(\|\widehat G_n-G\|<\lambda_{\min}(G)/2\). On this event, \(\widehat G_n\) is nonsingular,
\[
\|\widehat G_n^{-1}\|\leq \frac{2}{\lambda_{\min}(G)},
\qquad
\widehat G_n^{-1}-G^{-1}
=\widehat G_n^{-1}(G-\widehat G_n)G^{-1}
=O_P(\delta_n).
\]
Thus the definition of the loading tensor yields
\[
\begin{aligned}
\widehat\tau^{\mathrm{ex}}-G^{-1}b_n
&=\widehat G_n^{-1}(\widehat g_n-b_n)
  +(\widehat G_n^{-1}-G^{-1})b_n
=O_P(\delta_n),\\
G^{-1}b_n
&=\sum_{t=1}^{T}\sum_{\ell=0}^{p}
L_{\cdot,t,\ell}(\Pi)\overline\tau_{t,\ell,n}.
\end{aligned}
\]
This proves the first conclusion, including the stated rate. If
\(\tau_{i,t,\ell}=\tau_{i,-\ell}\) for every legal \(t\), the normalization in \(\text{thm:calendar-lag-loading}\) gives, for every \(j\in\{0,\ldots,p\}\),
\[
\begin{aligned}
e_j^{\top}G^{-1}b_n
&=\sum_{\ell=0}^{p}\left\{\sum_{t=1}^{T}L_{j,t,\ell}(\Pi)\right\}
  \left(n^{-1}\sum_i\tau_{i,-\ell}\right)\\
&=n^{-1}\sum_i\tau_{i,-j}.
\end{aligned}
\]
The coordinatewise homogeneous conclusion follows because \(d\) is fixed.

Finally, the exponent is sharp for the moment control available from the three published regularity conditions. Fix \(T=2\), \(p=0\), \(\mathcal S=\{(0,0),(0,1)\}\), and \(\Pi_{(0,0)}=\Pi_{(0,1)}=1/2\); direct substitution gives \(G(\Pi)=1/16>0\). Let \(B_n=\lceil n^q\rceil\) and partition the units into \(B_n\) blocks whose sizes \(m_{b,n}\) differ by at most one. Independently across blocks, draw a Rademacher variable \(V_b\) and a Bernoulli variable \(Z_b\) with success probability \(1/2\), independently of each other. For every unit in block \(b\), set \(W_i=(0,Z_b)\) and set
\[
Y_{i1}(w)=0,\qquad Y_{i2}(w)=V_b
\quad\text{for every }w\in\{0,1\}^2.
\]
Then \(\pi_i((0,0))=\pi_i((0,1))=1/2\) conditional on the full schedule, the second moments are bounded by one, and the mean model holds with \(\mu_{it}=0\) and \(\tau_{i,t,0}=0\). Within a block the observed schedule--treatment pairs are identical and nondegenerate, so \(\rho_{ij}=1\); across blocks they are independent, so \(\rho_{ij}=0\). Hence
\[
n^{-2}\sum_{i,j}\rho_{ij}
=n^{-2}\sum_{b=1}^{B_n}m_{b,n}^2
=\Theta(B_n^{-1})=\Theta(n^{-q}).
\]
The required weighted numerator moment
\(r_iD_{W_i}^{\top}JY_i=Z_bV_b/2\) is constant inside each block, centered, nondegenerate, and has sample-mean standard deviation of order \(n^{-q/2}\). Moreover, this fluctuation survives in the coefficient. Put
\[
\overline Z_n=n^{-1}\sum_iZ_i,\qquad
\overline V_n=n^{-1}\sum_iV_i,\qquad
S_n=n^{-1}\sum_i(Z_i-\tfrac12)V_i.
\]
Since \(\xi_b=(2Z_b-1)V_b\) are independent Rademacher variables,
\[
S_n=\frac{1}{2n}\sum_{b=1}^{B_n}m_{b,n}\xi_b,
\qquad
\sigma_n^2:=\mathbb E[S_n^2]
=\frac{1}{4n^2}\sum_{b=1}^{B_n}m_{b,n}^2
=\Theta(n^{-q}).
\]
Independence and centering across blocks give
\(\mathbb E[S_n^4]\leq3\sigma_n^4\). This follows by expanding the fourth power: only terms with indices occurring in pairs remain, and each Rademacher fourth moment equals the squared second moment. With
\(A_n=\{S_n^2\geq\sigma_n^2/2\}\), Cauchy--Schwarz gives
\[
\sigma_n^2=\mathbb E[S_n^2]
\leq \frac{\sigma_n^2}{2}
  +\{\mathbb E[S_n^4]\mathbb P_n(A_n)\}^{1/2},
\]
and consequently \(\mathbb P_n(A_n)\geq1/12\). Thus
\(S_n\) is not \(o_P(n^{-q/2})\), and it cannot be \(O_P(n^{-r})\) for any \(r>q/2\).

In this two-period design, let
\[
E_n=\{0<\overline Z_n<1\}.
\]
There are \(B_n\) independent block-level treatment draws and every block has positive size, so
\[
\mathbb P_n(E_n^c)=2^{1-B_n}\longrightarrow0.
\]
On \(E_n\), the sample normal matrix is nonsingular and the exact fixed-effect coefficient is the slope from the cross-unit regression of the outcome change \(V_i\) on the treatment change \(Z_i\):
\[
\widehat\tau^{\mathrm{ex}}
=\frac{n^{-1}\sum_i(Z_i-\overline Z_n)V_i}
       {n^{-1}\sum_i(Z_i-\overline Z_n)^2}.
\]
On \(E_n^c\), retain any coefficient selected by the definition's arbitrary-minimizer convention; because \(\mathbb P_n(E_n^c)\to0\), this selection does not alter any conclusion in probability below. The same block-variance calculation gives
\(\overline Z_n-\tfrac12=O_P(n^{-q/2})\) and
\(\overline V_n=O_P(n^{-q/2})\). Hence, on \(E_n\),
\[
n^{-1}\sum_i(Z_i-\overline Z_n)V_i
=S_n-(\overline Z_n-\tfrac12)\overline V_n
=S_n+O_P(n^{-q}),
\]
whereas the binary identity gives
\[
n^{-1}\sum_i(Z_i-\overline Z_n)^2
=\overline Z_n(1-\overline Z_n)
=\frac14+O_P(n^{-q}).
\]
The coefficient target is zero, so, on \(E_n\),
\(\widehat\tau^{\mathrm{ex}}=4S_n+o_P(n^{-q/2})\). Since \(\mathbb P_n(E_n^c)\to0\), the preceding fourth-moment lower bound for \(S_n\) implies that \(\widehat\tau^{\mathrm{ex}}\) is not \(o_P(n^{-q/2})\), irrespective of the arbitrary minimizer selected on \(E_n^c\). This block construction proves that no uniformly faster coefficient rate follows from the three published regularity conditions alone; it invokes no central limit theorem and makes no asymptotic-normality claim.
